%% =========================================================
%% Chapter: Entropy and energy transport on symmetric replicas
%% Source: OpenAI, "A two-dimensional area law from a global spectral gap"
%%         (September 24, 2026), Proposition 7.4, 06-transport.tex lines 377-779
%% =========================================================

\chapter{Entropy and Energy Transport on Symmetric Replicas}
\label{ch:replica_transport}

This chapter assembles the entropy and energy transport estimate
of~\cite[Proposition~7.4]{OpenAI2026AreaLaw} on the symmetric subspace
$\mathcal S_k$ of $k$ replicas, from the matrix estimates of
Section~\ref{sec:matrix_mean_trees}. The inputs of the mean trees are products
over bands of replica metrics, the moves transfer a subsystem $x$ of the middle
region $Y$ to one of the outer regions, and $\eta$ denotes the conditional mutual
information of a move evaluated at a one-copy vector $\theta$. For a density
matrix $\rho$ on $\mathcal S_k$, $\int f\,d\mu_\rho$ denotes the coherent
integral $D_k\int f(\theta)\operatorname{Tr}(\rho P_{\theta,k})\,d\theta$, where
$P_{\theta,k}$ projects onto $\theta^{\otimes k}$ and $D_k=\dim\mathcal S_k$.

The source proves the relative coherent pin and the skew bound in its replica
section. The assembled estimate is first proved assuming their conclusions in the
form used at a single band; both forms are then derived from the replica lemmas,
which gives the proposition without hypotheses beyond those of the source.

\section{The logarithmic passage}

\begin{lemma}[Operator Jensen inequality on symmetric replicas]
    \label{lem:replica_coherent_log_jensen}
    \lean{TensorPower.coherentAverage_log_le, TensorPower.trace_mul_coherentAverage,
      TensorPower.log_symDim_isBigO}
    \leanok
    \uses{thm:coherent_resolution}
    Let $\mathcal Q_k(f)=D_k\int f(\theta)P_{\theta,k}\,d\theta$ and let $\Pi$ be
    the projection onto $\mathcal S_k$. For continuous $f>0$,
    \[
      \mathcal Q_k(\log f)\le\Pi\log\bigl(\mathcal Q_k(f)+(I-\Pi)\bigr)\Pi .
    \]
    Moreover $\log D_k=O(\log(k+1))$.
\end{lemma}

\begin{proof}\leanok
    Write $\log Y=\int_0^\infty[(1+s)^{-1}-(Y+s)^{-1}]\,ds$. For each $s$ the
    positivity of $g(X-g^{-1})P_\theta(X-g^{-1})$ with
    $X=\Pi(Y+s)^{-1}\Pi$, averaged over $\theta$, gives the resolvent form of the
    inequality; Fubini's theorem exchanges the two integrals
    (\texttt{06-transport.tex} lines 526--546). The dimension of $\mathcal S_k$
    is at most $(k+1)^{\dim}$.
\end{proof}

\begin{lemma}[Logarithmic passage]\label{lem:replica_log_passage}
    \lean{TensorPower.realCoherentIntegral_sub_log_le_re_trace_mul_log}
    \leanok
    \uses{lem:replica_coherent_log_jensen}
    Let $X$ be positive definite and commute with copy permutations, and let
    $\varphi$ be continuous with $X\ge e^{\varphi(\theta)}P_{\theta,k}$ for every
    unit vector $\theta$. For every density matrix $\rho$ on $\mathcal S_k$,
    $\int\varphi\,d\mu_\rho-\log D_k\le\operatorname{Tr}(\rho\log X)$.
\end{lemma}

\begin{proof}\leanok
    Integrating the pin gives $D_k^{-1}\mathcal Q_k(e^\varphi)\le X$. Adding
    $\varepsilon(I-\Pi)$ with $\varepsilon I\le X$ keeps the inequality and makes
    the left side positive definite; operator monotonicity of the logarithm and
    the Jensen inequality conclude (\texttt{06-transport.tex} lines 520--575).
\end{proof}

\section{The assembled estimate}

\begin{lemma}[Band metrics on the symmetric subspace]\label{lem:replica_symmetric_metric}
    \lean{Matrix.mul_cfc_eq_mul_cfc_of_mul_eq, Matrix.cfc_proj_mul_add,
      Matrix.PosDef.proj_mul_add_one_sub, Matrix.proj_mul_add_one_sub_mul,
      Matrix.PosDef.whiten_proj_mul_add_one_sub,
      TensorPower.ReplicaTransport.symBandMetric,
      TensorPower.ReplicaTransport.posDef_symBandMetric,
      TensorPower.ReplicaTransport.commute_permOp_symBandMetric,
      TensorPower.ReplicaTransport.TransportData.crossBandCommute_iff,
      TensorPower.ReplicaTransport.le_whiten_symBandMetric,
      TensorPower.ReplicaTransport.mul_skewSquare_symBandMetric}
    \leanok
    Let $\Pi$ be the projection onto $\mathcal S_k$ and let $A$ be a band metric, which
    commutes with $\Pi$. The operator $\widetilde A=\Pi A+(I-\Pi)$ is positive definite,
    agrees with $A$ on $\mathcal S_k$ and commutes with the copy permutations. Two such
    operators $\widetilde A,\widetilde B$ commute exactly when $A$ and $B$ commute on
    $\mathcal S_k$. Moreover
    $\widetilde A^{-1/2}\widetilde B\widetilde A^{-1/2}=\Pi A^{-1/2}BA^{-1/2}+(I-\Pi)$, so
    every lower bound $Q\le A^{-1/2}BA^{-1/2}$ with $\Pi Q\Pi=Q$ holds for the extended
    operators, and for every $X$ commuting with $\Pi$ the skew squares of $\widetilde A$ and
    of $A$ have the same expectation in every density matrix on $\mathcal S_k$.
\end{lemma}

\begin{proof}\leanok
    If $P$ commutes with a Hermitian $Z$ and $PY=PZ$ for a Hermitian $Y$, then $PY^k=PZ^k$
    for all $k$; interpolating $f$ by one polynomial on the two finite spectra gives
    $Pf(Y)=Pf(Z)$. Applied with $P=\Pi$ and $P=I-\Pi$ this computes
    $f(\Pi Y+(I-\Pi)Z)=\Pi f(Y)+(I-\Pi)f(Z)$, in particular the powers of $\widetilde A$.
    The identity $\widetilde A\widetilde B=\Pi AB+(I-\Pi)$ gives the commutation criterion,
    and positivity follows by splitting a vector into its components in $\mathcal S_k$ and
    its complement.
\end{proof}

\begin{lemma}[Transport states on replicas]\label{lem:replica_transport_states}
    \lean{TensorPower.ReplicaTransport.TransportData.posSemidef_state,
      TensorPower.ReplicaTransport.TransportData.trace_state,
      TensorPower.ReplicaTransport.TransportData.symProj_mul_state,
      TensorPower.ReplicaTransport.TransportData.hasDerivAt_exactDerivative}
    \leanok
    \uses{lem:transport_exact_derivative, lem:transport_states_commutant}
    For admissible transport data the transport states are density matrices
    supported on $\mathcal S_k$, and the exact derivative formula of
    Lemma~\ref{lem:transport_exact_derivative} holds for the replica inputs.
\end{lemma}

\begin{proof}\leanok
    All inputs are positive definite and commute with copy permutations, so the
    symmetrizer fixes the transport states by
    Lemma~\ref{lem:transport_states_commutant}.
\end{proof}

\begin{theorem}[Entropy and energy transport, assuming Lemmas 6.3 and 6.4]
    \label{thm:replica_transport_conditional}
    \lean{TensorPower.ReplicaTransport.transport_of_relativePin_of_splitSkewBound,
      TensorPower.ReplicaTransport.TransportData.exists_entropyGain_le_exactDerivative_of_relativePin,
      TensorPower.ReplicaTransport.TransportData.exists_energy_le_of_splitSkewBound,
      TensorPower.ReplicaTransport.TransportData.smul_coherentProj_le_bandRelRatio,
      TensorPower.ReplicaTransport.TransportData.cfc_log_relRatio_eq_sum}
    \leanok
    \uses{lem:replica_log_passage, lem:replica_transport_states,
      lem:transport_energy_block, lem:replica_symmetric_metric}
    Assume the conclusion of the relative coherent pin~\cite[Lemma~6.3]{OpenAI2026AreaLaw}
    and of the skew bound~\cite[Lemma~6.4]{OpenAI2026AreaLaw} at a split leaf, and that
    metrics from different bands commute on $\mathcal S_k$, for every $k$, as in the
    source. The mean trees are evaluated on the band metrics extended by the identity off
    $\mathcal S_k$ (Lemma~\ref{lem:replica_symmetric_metric}), which agree with the
    source's operators on $\mathcal S_k$; an earlier formalization assumed commutation on
    the full tensor power instead~\cite{gap:oai26_transport_crossband_commutation}.
    Then the conclusions of~\cite[Proposition~7.4]{OpenAI2026AreaLaw} hold: with
    $a=2t$, $\ell\ge1$ bounding $\log(e\dim x)$ for every transferred subsystem and
    $\log(e\dim D_i)$ for every split term, and $a\ell$ below a universal
    threshold, for every $k$, nonzero $x\in\mathcal S_k$ and $0<p<1$,
    \[
      -\partial_p\log N(p)^2\ge ka\sum_{h,g}w_h\int m_{1/4}(u)\int\sum_cq_{c|h}
      \eta_{h,c,g}\,d\mu_{\sigma_{h,u}}\,du-CkaKa^{1/4}\ell^C-\beta_k
    \]
    with $\beta_k=O(\log(k+1))$ independent of $p$ and $x$; and if
    $\bar Hx=E_0x$, then
    $\langle v,\bar Hv\rangle\le2E_0+Ca^2\ell^C\sum_iW_i(p)\sum_{j}\pi_j\int
    m_{1/4}\int\eta_{i,j}^{1/8}\,d\mu_{\sigma_{j,u}}\,du+o_k(1)$, uniformly in $p$
    and $x$.
\end{theorem}

\begin{proof}\leanok
    The extended band metrics of different bands commute, so the relative ratio factors
    over bands, $\log C_h=\sum_g\log C_{h,g}$. The
    relative pin and projection transfer give the conditional pin of each band
    (\texttt{06-transport.tex} lines 513--531), and the logarithmic passage turns
    it into the entropy lower bound (lines 495--583). For the energy, the
    support of each term commutes with every unsplit band, a single band remains
    at a split leaf (lines 732--735), the skew bound controls its symbol cost
    (lines 736--748), and Lemma~\ref{lem:transport_energy_block} applies term by
    term with $\bar Hx=E_0x$ (lines 588--766).
\end{proof}

\begin{theorem}[Entropy and energy transport, assuming the split-leaf skew bound]
    \label{thm:replica_transport_skew_conditional}
    \lean{Matrix.smul_vecMulVec_le_of_forall,
      TensorPower.ReplicaTransport.fourPartition_of_isPartition,
      TensorPower.ReplicaTransport.isBigO_neg_log_inv_poly,
      TensorPower.ReplicaTransport.relativePinBound,
      TensorPower.ReplicaTransport.transport_of_splitSkewBound}
    \leanok
    \uses{thm:replica_transport_conditional, thm:relative_pin, lem:movement_eta_regional}
    The relative coherent pin holds in the form assumed in
    Theorem~\ref{thm:replica_transport_conditional}. Hence the conclusions
    of~\cite[Proposition~7.4]{OpenAI2026AreaLaw} hold assuming only the skew bound at a
    split leaf.
\end{theorem}

\begin{proof}\leanok
    For a move of $x\subseteq Y$ to $P$, the parts $P,x,Y\setminus x,F$ partition the
    sites, and Theorem~\ref{thm:relative_pin} bounds the relative metric with
    $\eta_\theta=S_\theta(x|P)+S_\theta(x|Y\setminus x)=I_\theta(x:F|P)$, the entropy of the
    move (\texttt{06-transport.tex} lines 315--328). A move to $F$ is the same after
    exchanging the outer parts, and a choice making no move has relative metric $I$, which
    dominates every coherent projector. A quadratic-form bound
    $c\,|\langle\theta^{\otimes k},w\rangle|^2\le\langle w,Mw\rangle$ for Hermitian $M$ is the
    operator bound $cP_{\theta,k}\le M$. The polynomial factor $p(k)$ gives
    $-\log p(k)^{-1}=O(\log(k+1))$, and $\log(e\dim x)$ is the logarithm of $e$ times the
    product of the site dimensions.
\end{proof}

\begin{corollary}[Integrated form]\label{cor:replica_transport_integrated}
    \lean{TensorPower.ReplicaTransport.TransportData.integral_le_log_filteredNormSq_sub}
    \leanok
    \uses{lem:transport_integrated_derivative}
    For $0\le p_0\le p_1\le1$ and every integrable function $L$ on $[0,1]$ with
    $L(p)\le-\partial_p\log N(p)^2$ for $0<p<1$,
    $\int_{p_0}^{p_1}L\le\log N(p_0)^2-\log N(p_1)^2$.
\end{corollary}

\begin{proof}\leanok
    Integrate the pointwise bound and apply
    Lemma~\ref{lem:transport_integrated_derivative} (\texttt{06-transport.tex}
    lines 427--429 and 769--779).
\end{proof}

\begin{lemma}[Integrability of the entropy gain]\label{lem:replica_entropy_gain_integrable}
    \lean{Matrix.MeanTree.hasDerivAt_eval_update_path,
      Matrix.Transport.hasDerivAt_neg_log_filteredNormSq_leafPerturbation,
      TensorPower.ReplicaTransport.TransportData.aestronglyMeasurable_integral_state,
      TensorPower.ReplicaTransport.TransportData.abs_integral_state_le,
      TensorPower.ReplicaTransport.TransportData.intervalIntegrable_entropyGain}
    \leanok
    \uses{lem:replica_transport_states}
    For fixed $k$ and nonzero $x\in\mathcal S_k$, the entropy-gain term of
    \cite[Proposition~7.4]{OpenAI2026AreaLaw} is integrable in $p$ on $[0,1]$.
\end{lemma}

\begin{proof}\leanok
    Each summand is $G(p)=\int m_{1/4}(u)\operatorname{Tr}(\sigma_{h,u}Q)\,du$ with $Q$
    Hermitian. Replacing the old input $A_h$ by $A_h^{1/2}e^{\delta Q}A_h^{1/2}$ changes
    $-\log N(p)^2$ at $\delta=0$ at rate $(1-p)w_hG(p)$, by the norm-derivative formula and
    the leaf identity. For fixed $\delta$ the perturbed norm is continuous in $p$, so
    $(1-p)w_hG(p)$ is a pointwise limit of continuous functions on $(0,1)$ and hence
    measurable. The states are density matrices, so $|G(p)|\le\|Q\|/2$. The source asserts
    the integration (\texttt{06-transport.tex} lines 427--429) and bounds the integrand
    (lines 769--779); measurability is the remaining point.
\end{proof}

\begin{theorem}[The skew bound at a split leaf]\label{thm:replica_split_skew_bound}
    \lean{TensorPower.re_coherentIntegral_ofReal,
      TensorPower.mul_skewOperator_eq_mul_skewSquare,
      TensorPower.ReplicaTransport.posSemidef_of_localLift,
      TensorPower.ReplicaTransport.exists_tendsto_zero_forall_le,
      TensorPower.ReplicaTransport.eventually_re_trace_skewSquare_le,
      TensorPower.ReplicaTransport.splitSkewBound}
    \leanok
    \uses{thm:replica_transport_skew_conditional}
    Let $0\le h\le1$ have designated support $D$ meeting $Y$ and exactly one of the outer
    parts of a partition, let $0<t<1/4$, $a=2t$ and $a\log(e\dim D)\le1/8$. Then, uniformly
    over density matrices $\rho$ on $\mathcal S_k$,
    $\operatorname{Tr}(\rho\mathsf D_k)\le 2C a^2\log^4(e\dim D)\int\eta^{1/8}\,d\mu_\rho+r_k$
    with $r_k\to0$, where $\mathsf D_k$ is the skew square of the band metric and $\eta$ is
    the entropy of the split.
\end{theorem}

\begin{proof}\leanok
    For a split towards $P$ put $P_0=D\cap P$ and $x=D\cap Y$. Then $D=P_0\cup x$, so $h$ is
    the lift of an operator $0\le K\le1$ on $P_0x$ and $\dim D=(\dim P_0)(\dim x)$. The skew
    operator of~\cite[Lemma~6.4]{OpenAI2026AreaLaw} takes its absolute values on
    $\mathcal S_k$; for $Z$ commuting with the projection $\Pi$ onto $\mathcal S_k$,
    $\Pi|\Pi Z|=\Pi|Z|$, so the two skew operators have the same expectation in every
    density matrix on $\mathcal S_k$. The entropy $S_\theta(x|P)+S_\theta(x|Y\setminus x)$
    of the skew bound is $I_\theta(x:F|P)$. A split towards $F$ is the same with the outer
    parts exchanged. The skew bound holds with an additive $\varepsilon$ for all large $k$;
    for fixed $k$ the excess over symmetric density matrices is bounded by the operator norm
    of the skew square, and its supremum is the remainder $r_k$.
\end{proof}

\begin{theorem}[Entropy and energy transport]\label{thm:replica_transport}
    \lean{TensorPower.ReplicaTransport.transport}
    \leanok
    \uses{thm:replica_transport_skew_conditional, thm:replica_split_skew_bound,
      cor:replica_transport_integrated, lem:replica_entropy_gain_integrable}
    The conclusions of~\cite[Proposition~7.4]{OpenAI2026AreaLaw} hold under the hypotheses
    of the source: admissible finite history and choice data, cross-band commutation on
    $\mathcal S_k$, positive contractions with designated supports compatible with the
    terminal partitions, and $\ell\ge1$ bounding the logarithmic dimensions with $a\ell$
    below a universal threshold. Namely, for every $k$ and nonzero $x\in\mathcal S_k$: at
    every $0<p<1$ the exact derivative, the entropy-gain lower bound and the energy bound
    hold, and the entropy-gain lower bound integrates over every $[p_0,p_1]\subseteq[0,1]$
    to $\log N(p_0)^2-\log N(p_1)^2$.
\end{theorem}

\begin{proof}\leanok
    Combine Theorem~\ref{thm:replica_transport_skew_conditional} with
    Theorem~\ref{thm:replica_split_skew_bound}. For the integrated form, apply
    Corollary~\ref{cor:replica_transport_integrated} with the entropy-gain lower bound,
    which is integrable by Lemma~\ref{lem:replica_entropy_gain_integrable}.
\end{proof}
